\documentclass[withoutpreface]{cumcmthesis}
\usepackage[framemethod=TikZ]{mdframed}
\usepackage{url}
\usepackage{subcaption}

% 字体设置（使用系统自带字体）
\setCJKmainfont{SimSun}[Path=./]
\setCJKsansfont{SimHei}[Path=./]
\setCJKmonofont{FangSong}[Path=./]

\tihao{B}
\baominghao{2025001}
\schoolname{某大学}
\membera{队员A}
\memberb{队员B}
\memberc{队员C}
\supervisor{指导教师}
\yearinput{2025}
\monthinput{09}
\dayinput{15}

\begin{document}

\maketitle

\begin{abstract}
本文针对...问题进行了系统研究。首先对问题进行了深入分析，建立了相应的数学模型...（摘要正文300-500字）

\textbf{关键词}: 关键词1；关键词2；关键词3；关键词4；关键词5
\end{abstract}

\section{1 问题重述}

\subsection{1.1 研究背景}
\ldots

\subsection{1.2 问题描述}
\ldots

\section{2 问题分析}

\subsection{2.1 总体分析}
\ldots

\subsection{2.2 求解思路}
\ldots

\section{3 模型假设与符号说明}

\subsection{3.1 模型假设}
\begin{enumerate}
\item 假设1：\ldots
\item 假设2：\ldots
\end{enumerate}

\subsection{3.2 符号说明}
\begin{table}[!htbp]
\centering
\caption{主要符号说明}
\begin{tabular}{ccp{8cm}}
\toprule
符号 & 意义 & 说明 \\
\midrule
$d$ & 外延层厚度 & 单位：cm \\
$\Delta\sigma$ & 相邻干涉峰波数间距 & 单位：cm$^{-1}$ \\
$n$ & 外延层折射率 & 由Sellmeier色散模型计算 \\
\bottomrule
\end{tabular}
\end{table}

\section{4 模型建立}

\subsection{4.1 模型名称}
\ldots

\section{5 模型求解}

\subsection{5.1 算法设计}
\ldots

\section{6 计算结果}

\subsection{6.1 主要结果}
\ldots

\section{7 可靠性分析}

\section{8 结论}

\section{9 参考文献}

\begin{thebibliography}{9}
\addcontentsline{toc}{section}{参考文献}
\end{thebibliography}

\newpage

\begin{appendices}

\section{Python求解代码}

\begin{lstlisting}[language=python]
# 代码内容
\end{lstlisting}

\end{appendices}

\end{document}